% GENERATED, edit the French unit or the translation table
\mxunit{seconde-s03}{Straight lines in the plane}\mxsubtitle{1--16 October 2026}
\section{Direction vector and Cartesian equation}
\begin{keypoint}[Definition]
A direction vector of a line is a \textbf{nonzero} vector having the direction of that line. If $A\ne B$, $\vect{AB}$ directs $(AB)$.
Every line has an equation $ax+by+c=0$ with $\pt{a;b}\ne\pt{0;0}$. A direction vector is $\vc{-b}{a}$.
\end{keypoint}
\begin{example}For $A\pt{-2;4}$ and $B\pt{6;2}$, calculate $\vect{AB}$ then check that $\vec u\pt{4;-1}$ directs $(AB)$.\end{example}
\mxanswerspace[2.4]{$\vect{AB}=\pt{8;-2}=2\vec u$. The vector $\vec u$ is nonzero and collinear with $\vect{AB}$; it therefore directs $(AB)$.}
\begin{keypoint}[Establishing an equation with the determinant]
The line through $A\pt{x_A;y_A}$ with direction vector $\vec u\pt{\alpha;\beta}\ne\vec0$ has equation
\[\beta(x-x_A)-\alpha(y-y_A)=0.\]
\end{keypoint}
\mxwriteproof[4]{A point $M\pt{x;y}$ belongs to the line if and only if $\vect{AM}$ is collinear with $\vec u$, including for $M=A$. Now
\[\det(\vect{AM},\vec u)=(x-x_A)\beta-(y-y_A)\alpha.\]
Setting the determinant to zero gives the stated equation. Expanding, $a=\beta$, $b=-\alpha$ and $c=-\beta x_A+\alpha y_A$; since $\vec u\ne\vec0$, $\pt{a;b}\ne\pt{0;0}$.}
\begin{pitfall}The zero vector directs no line. The sign of the determinant depends on the order of the vectors, but whether it vanishes does not.\end{pitfall}

\mxsection{Determining the equation of a line}
\begin{method}
From a point and a direction vector: choose $a,b$ then replace $x,y$ by the coordinates of the point to find $c$. From two distinct points: first calculate a direction vector.
\end{method}
\begin{example}Determine a Cartesian equation of the line through $A\pt{2;-1}$ with direction vector $\vec u\pt{3;2}$.\end{example}
\mxanswerspace[4]{$a=2$, $b=-3$ will do. The line therefore has an equation $2x-3y+c=0$. Since $A$ belongs to it, $4+3+c=0$, so $c=-7$. An equation is $2x-3y-7=0$. Check: $2\times2-3\times(-1)-7=0$ and $\pt{3;2}$ is indeed a direction vector.}
\begin{example}Determine a Cartesian equation of $(AB)$ for $A\pt{3;3}$ and $B\pt{2;-2}$.\end{example}
\mxanswerspace[4]{$\vect{AB}=\pt{-1;-5}$ directs the line. We can choose $\pt{a;b}=\pt{5;-1}$. With $A$, $15-3+c=0$, so $c=-12$. Thus $5x-y-12=0$. For $B$, $10+2-12=0$, which confirms the result.}
\begin{example}Does the point $P\pt{1;2}$ belong to $d:2x-3y+4=0$? Does the point $Q\pt{2;1}$ belong to it?\end{example}
\mxanswerspace[2.4]{For $P$, $2-6+4=0$: yes. For $Q$, $4-3+4=5\ne0$: no.}

\mxsection{Slope-intercept equation, slope and affine function}
\begin{keypoint}[Definitions]
A non-vertical line has a unique slope-intercept equation $y=mx+p$. $m$ is the slope, $p$ the y-intercept. It represents the affine function $x\mapsto mx+p$; $\pt{1;m}$ is a direction vector.
A vertical line has an equation $x=k$; it is not the graph of a function of $x$. A horizontal line has an equation $y=k$.
\end{keypoint}
\begin{keypoint}[Calculating a slope]
If $x_A\ne x_B$, the slope of $(AB)$ is $m=\dfrac{y_B-y_A}{x_B-x_A}$. To find $p$, we then use $y_A=mx_A+p$.
\end{keypoint}
\begin{example}Determine the slope-intercept equation of $2x+3y-1=0$, then draw $d:y=-2x+3$.\end{example}
\mxanswerspace[2.4]{$3y=-2x+1$, so $y=-\frac23x+\frac13$. For $d$, plot $\pt{0;3}$ and $\pt{1;1}$, then draw the line; a move of $1$ along the x-axis corresponds to $-2$ on the y-axis.}
\begin{center}\begin{tikzpicture}\begin{axis}[mxgridaxis,xmin=-1,xmax=4,ymin=-4,ymax=4,width=10cm,height=6cm,xtick={-1,...,4},ytick={-4,...,4}]
\ifmxteacher\addplot[thick,domain=-.5:3.5]{-2*x+3};\addplot[only marks]coordinates{(0,3)(1,1)};\fi
\end{axis}\end{tikzpicture}\end{center}
\begin{example}Find the affine function $f$ such that $f(0)=3$ and $f(-1)=5$, then the affine function $g$ whose graph passes through $C\pt{-1;2}$ and $E\pt{5;-2}$.\end{example}
\mxanswerspace[2.4]{$f(x)=-2x+3$, since $m=(3-5)/(0+1)=-2$. For $g$, $m=(-2-2)/(5+1)=-2/3$ and $p=2-2/3=4/3$, so $g(x)=-\frac23x+\frac43$.}

\mxsection{Relative positions of two lines}
\begin{keypoint}[Properties]
Two lines are either intersecting, strictly parallel or coincident. For $y=mx+p$ and $y=m'x+p'$:
\begin{itemize}\item if $m\ne m'$, they intersect;\item if $m=m'$ and $p\ne p'$, they are strictly parallel;\item if $m=m'$ and $p=p'$, they coincide.\end{itemize}
For Cartesian equations, we compare the directions with the determinant. If the directions are collinear, we then test the membership of a point to distinguish coincident lines from strictly parallel ones.
\end{keypoint}
\begin{example}Determine an equation of $d$ through $A\pt{2;3}$ and parallel to $d':x-3y-3=0$.\end{example}
\mxanswerspace[3.2]{$d'$ has direction vector $\pt{3;1}$. We look for $d:x-3y+c=0$. Since $A\in d$, $2-9+c=0$, so $c=7$. Thus $d:x-3y+7=0$. The two lines are strictly parallel, because the constants differ for the same coefficients of $x,y$.}
\begin{example}Study the relative positions:
\begin{enumerate}\item $2x-y+1=0$ and $4x-2y+2=0$;\item $x=2$ and $x=-1$;\item $y=2x-3$ and $y=-x+3$.\end{enumerate}\end{example}
\mxanswerspace[4]{1. The second equation is twice the first: coincident lines. 2. Two distinct vertical lines: strictly parallel. 3. The slopes $2$ and $-1$ differ: intersecting. Their intersection satisfies $2x-3=-x+3$, so $x=2$ and $y=1$.}

\mxsection{Intersection: solving by substitution}
\begin{keypoint}[Definition]
Solving a system of two equations in $x,y$ means finding all the pairs $\pt{x;y}$ satisfying both equations simultaneously. For two equations of lines: intersecting lines give one solution pair; strictly parallel lines give no solution; coincident lines give infinitely many solutions.
\end{keypoint}
\begin{method}[Substitution]
Isolate one unknown in one equation, substitute its expression into the other, solve, then find the second unknown. Check the pair in both original equations.
\end{method}
\begin{example}Solve $\begin{cases}-2x+y+3=0,\\3x-5y-2=0.\end{cases}$\end{example}
\mxanswerspace[4]{The first equation gives $y=2x-3$. Substituting into the second: $3x-5(2x-3)-2=0$, that is $-7x+13=0$. Thus $x=13/7$ then $y=5/7$. The solution set is $S=\{\pt{13/7;5/7}\}$. The substitutions give $(-26+5+21)/7=0$ and $(39-25-14)/7=0$.}
\begin{example}Solve by substitution:
\[\text{(a) }\begin{cases}-2x+y+3=0\\x-4y+5=0\end{cases}\qquad
\text{(b) }\begin{cases}2x+3y-10=0\\4x+5y-16=0.\end{cases}\]\end{example}
\mxanswerspace[4]{(a) $y=2x-3$ then $x-4(2x-3)+5=0$, so $x=17/7$, $y=13/7$.\\(b) $x=(10-3y)/2$; $2(10-3y)+5y-16=0$, so $y=4$, then $x=-1$.}

\mxsection{Linear combination and Python}
\begin{method}[Linear combination]
Multiply the equations by nonzero numbers to obtain opposite coefficients for one unknown; add them to eliminate that unknown. Keep one equation of the system to calculate the other unknown.
\end{method}
\begin{example}Solve $\begin{cases}2x+3y+1=0\\3x-4y-2=0.\end{cases}$ Then redo the two systems of the previous example by linear combination.\end{example}
\mxanswerspace[4.8]{Multiplying the first equation by $3$ and the second by $-2$ gives $6x+9y+3=0$ and $-6x+8y+4=0$. The sum gives $17y+7=0$, so $y=-7/17$. Then $2x-21/17+1=0$, so $x=2/17$.
For (a), add four times the first equation to the second: $-7x+17=0$, so $\pt{x;y}=\pt{17/7;13/7}$. For (b), subtract twice the first from the second: $-y+4=0$, so $\pt{x;y}=\pt{-1;4}$.}
\begin{lstlisting}
def determinant(u, v):
    return u[0]*v[1] - u[1]*v[0]

def alignes(A, B, C):
    AB = (B[0]-A[0], B[1]-A[1])
    AC = (C[0]-A[0], C[1]-A[1])
    return determinant(AB, AC) == 0
\end{lstlisting}
\begin{example}Predict \texttt{alignes((0,0), (2,1), (6,3))} then \texttt{alignes((0,0), (2,1), (6,4))}.\end{example}
\mxanswerspace[2]{\texttt{True}, then \texttt{False}: the determinants are $0$ and $2$. The exact test is suitable here for integer coordinates; floating-point approximations require taking numerical errors into account.}
